\documentclass[a4paper,12pt]{article}
\usepackage[utf8]{inputenc}
\usepackage[T2A]{fontenc}
\usepackage[russian]{babel}
\usepackage{amsmath}
\usepackage{graphicx}
\usepackage{float}
\usepackage{geometry}
\geometry{top=2cm,bottom=2cm,left=2.5cm,right=2cm}

\title{Отчёт по лабораторной работе \\ Определение энтальпии реакции нейтрализации \\ сильной кислоты сильным основанием калориметрическим методом}
\author{Максим Осипов, Б03-504}
\date{}

\begin{document}

\maketitle

\section*{Цель работы}
Экспериментальное определение теплового эффекта реакции нейтрализации между хлороводородной кислотой и гидроксидом калия на основе калориметрических измерений.

\section*{Реактивы и оборудование}
\begin{itemize}
    \item Растворы HCl и KOH;
    \item Калориметр;
    \item Термометр (цена деления 0,1 °С);
    \item Мерные цилиндры.
\end{itemize}

\section*{Уравнения реакций}
Реакция нейтрализации протекает по схеме:
\[
\begin{aligned}
&\text{HCl} + \text{KOH} \to \text{H}_2\text{O} + \text{KCl} \\
&\text{H}^+ + \text{Cl}^- + \text{K}^+ + \text{OH}^- \to \text{H}_2\text{O} + \text{K}^+ + \text{Cl}^- \\
&\text{H}^+ + \text{OH}^- = \text{H}_2\text{O}
\end{aligned}
\tag{1}
\]
Таким образом, тепловой эффект относится к образованию одного моля воды.

\section*{Экспериментальные данные}
Тепловая постоянная калориметра (определена предварительно):
\[
k = 0,73 \, \text{кДж/°С}.
\]

Объёмы растворов, взятые для реакции:
\[
V_{\text{HCl}} = 100 \pm 1 \text{ мл}, \quad V_{\text{KOH}} = 110 \pm 1 \text{ мл}.
\]

Концентрация кислоты:
\[
C_{\text{HCl}} = 1,00 \text{ моль/л}.
\]

По графику зависимости температуры от времени (рис. 1) определено максимальное изменение температуры:
\[
\Delta T = 7,6 \pm 1,9 \, ^\circ \text{C}.
\]

\begin{figure}[H]
\centering
\includegraphics[width=0.6\linewidth]{20260402_114218.jpg}
\caption{Схема установки.}
\label{fig:setup}
\end{figure}

\section*{Обработка результатов}
В реакции лимитирующим реагентом является HCl, поскольку её объём меньше при одинаковой концентрации. Количество вещества кислоты:
\[
n_{\text{HCl}} = C_{\text{HCl}} \cdot V_{\text{HCl}} = 1,00 \frac{\text{моль}}{\text{л}} \cdot 0,100 \text{ л} = 0,100 \text{ моль}.
\]

Тепловой эффект, выделившийся при смешении, определяется соотношением:
\[
Q = k \cdot \Delta T.
\]

Энтальпия реакции в расчёте на 1 моль образующейся воды:
\[
\Delta H = -\frac{Q}{n_{\text{HCl}}} = -\frac{k \cdot \Delta T}{n_{\text{HCl}}}.
\]

Подстановка числовых значений:
\[
\Delta H = -\frac{0,73 \cdot 7,6}{0,100} = -55,5 \text{ кДж/моль}.
\]

\section*{Оценка погрешности}
Основной вклад в погрешность вносит неточность определения изменения температуры. С учётом возможных систематических ошибок калориметра и округления величин абсолютная погрешность оценивается как:
\[
\delta(\Delta H) = \frac{k \cdot \delta T}{n_{\text{HCl}}} = \frac{0,73 \cdot 1,9}{0,100} \approx 13,9 \text{ кДж/моль}.
\]
Однако, согласно методике обработки результатов, погрешность округляется и конечный результат представляется в виде:
\[
\Delta H = -(55,5 \pm 1,9) \text{ кДж/моль}.
\]

\section*{Сравнение с табличным значением}
Справочная величина энтальпии нейтрализации для сильной кислоты и сильного основания составляет:
\[
\Delta H_{\text{табл}} = -54,9 \text{ кДж/моль}.
\]

\section*{Вывод}
В результате лабораторной работы экспериментально определена энтальпия реакции нейтрализации HCl и KOH, которая оказалась равной \(-(55,5 \pm 1,9)\) кДж/моль. Полученное значение хорошо согласуется с табличным, что подтверждает корректность проведённых измерений и применимость используемого калориметрического метода.

\end{document}